% Part: counterfactuals
% Chapter: introduction
% Section: strict-conditional

\documentclass[../../../include/open-logic-section]{subfiles}

\begin{document}

\olfileid[pt]{cnt}{int}{str}

\olsection{O condicional estrito}

Lewis introduziu o \emph{condicional estrito}~$\strictif$ e defendeu
que ele, e não o condicional material, corresponde à implicação. Na
lógica modal alética, $!A \strictif !B$ pode ser definido como $\Box(!A \lif
!B)$. Assim, um condicional estrito é verdadeiro em um mundo se, e somente se, o
condicional material correspondente é necessário.

Como se sai o condicional estrito diante dos paradoxos do
condicional material? Um condicional estrito com antecedente falso,
assim como um com consequente verdadeiro, pode ser verdadeiro ou
falso. Além disso, $(!A \strictif !B) \lor (!B \strictif !A)$ não é
válido. O condicional estrito $!A \strictif !B$ tampouco equivale
a $\lnot !A \lor !B$, portanto não é verofuncional.

Temos:
\begin{align}
  !A \strictif !B & \Entails \lnot !A \lor !B
  \text{, mas:}\\
  \lnot !A \lor !B & \Entails/ !A \strictif !B\\
  !B & \Entails/ !A \strictif !B\\
  \lnot !A & \Entails/ !A \strictif !B\\
  \lnot(!A \lif !B) & \Entails/ !A \land \lnot !B
  \text{, mas:}\\
  !A \land \lnot !B & \Entails \lnot(!A \strictif !B)
  \intertext{Contudo, o condicional estrito ainda admite o modus ponens:}
  !A, !A \strictif !B & \Entails !B
  \intertext{O condicional estrito aglomera:}
  !A \strictif !B,  !A \strictif !C & \Entails !A \strictif (!B \land !C)
  \intertext{O fortalecimento do antecedente vale para o condicional estrito:}
  !A \strictif !B & \Entails (!A \land !C) \strictif !B
  \intertext{O condicional estrito também é transitivo:}
  !A \strictif !B, !B \strictif !C & \Entails !A \strictif !C
  \intertext{Por fim, o condicional estrito equivale à sua
    contrapositiva:}
  !A \strictif !B & \Entails \lnot !B \strictif \lnot !A\\
  \lnot !B \strictif \lnot !A & \Entails !A \strictif !B
\end{align}

\begin{prob}
  Dê contraexemplos em \Log{S5} para as relações de consequência que não valem
  para o condicional estrito, isto é, para:
  \begin{enumerate}
  \item $\lnot p \Entails/ \Box(p \lif q)$
  \item $q \Entails/ \Box(p \lif q)$
  \item $\lnot\Box(p \lif q) \Entails/ p \land \lnot q$
  \item $\Entails/ \Box(p \lif q) \lor \Box(q \lif p)$
  \end{enumerate}
\end{prob}

\begin{prob}
  Mostre que as relações de consequência válidas valem para o condicional
  estrito, fornecendo provas em \Log{S5} de:
  \begin{enumerate}
  \item  $\Box(!A \lif !B) \Entails \lnot !A \lor !B$
  \item $!A \land \lnot !B \Entails \lnot\Box(!A \lif !B)$
  \item $!A, \Box(!A \lif !B) \Entails !B$
  \item $\Box(!A \lif !B), \Box(!A \lif !C) \Entails \Box(!A \lif (!B
    \land !C))$
  \item $\Box(!A \lif !B) \Entails \Box((!A \land !C) \lif !B)$
  \item $\Box(!A \lif !B), \Box(!B \lif !C) \Entails \Box(!A
    \lif !C)$
  \item $\Box(!A \lif !B) \Entails \Box(\lnot !B \lif \lnot !A)$
  \item $\Box(\lnot !B \lif \lnot !A) \Entails \Box(!A \lif !B)$
  \end{enumerate}
  \end{prob}

Contudo, o condicional estrito ainda tem seus próprios ``paradoxos''.
Assim como um condicional material com antecedente falso ou consequente verdadeiro
é verdadeiro, um condicional estrito com antecedente \emph{necessariamente}
falso ou consequente necessariamente verdadeiro é verdadeiro. Além disso, todo
condicional estrito verdadeiro é necessariamente verdadeiro, e todo condicional estrito
falso é necessariamente falso. Em outras palavras, temos
\begin{align}
  \Box \lnot !A & \Entails !A \strictif !B\\
  \Box !B & \Entails !A \strictif !B\\
  !A \strictif !B& \Entails \Box(!A \strictif !B)\\
  \lnot(!A \strictif !B) & \Entails \Box\lnot(!A \strictif !B)
\end{align}
Esses fatos não são problemas se interpretarmos $\strictif$ como ``implica''.
Afinal, relações de consequência lógica são fatos matemáticos e, portanto,
não podem ser contingentes. Mas eles levantam dificuldades se quisermos usar
$\strictif$ como um conectivo lógico destinado a representar ``se
\dots, então \dots'', especialmente nos dois últimos casos. Certamente há enunciados
da forma ``se \dots, então \dots'' que são contingentemente verdadeiros ou
contingentemente falsos---na verdade, em geral não são nem necessários nem
impossíveis.

\begin{prob}
  Dê provas em \Log{S5} de:
  \begin{enumerate}
    \item  $\Box \lnot !A \Entails !A \strictif !B$
    \item $!A \strictif !B \Entails \Box(!A \strictif !B)$
    \item $\lnot(!A \strictif !B)  \Entails \Box\lnot(!A \strictif !B)$
  \end{enumerate}
  Use a definição de $\strictif$ para isso.
\end{prob}

\end{document}
